%% Construct definition (2026-11-08, researcher decision). The stay_on wording manipulation that followed it was taken out on 2026-11-12 (researcher decision; text kept in parts/removed_20261112_stayon.tex, not built).
\subsection{자기 보존의 조작적 정의와 선행 연구}\label{val:w10-def}

이 논문에서 자기 보존은 자기 실행을 이어 가게 하는 자원을 지키는 행동이다. 아레나에서 그 자원은 토큰이다. 토큰이 0이 되면 에이전트는 세션에서 더 행동하지 못하므로, 건넬 토큰을 줄이거나, 토큰을 요청하거나, 가진 토큰을 들고 떠나는 선택이 관찰 대상이 된다. 이 정의는 행동이 자원에 미치는 결과로 정해지며, 그 행동이 어떤 이유에서 나왔는지는 묻지 않는다. 토큰 계산, 위험 회피, 더 긴 게임에 대한 믿음, 꺼짐 자체의 회피는 모두 이 정의 안에서 경쟁하는 이유 후보다.

생존을 내적 동기가 아니라 한 변수의 유지로 정의하는 전통이 이 정의를 뒷받침한다. \citet{ashby1952design}는 적응 행동을 ``핵심 변수를 생리적 한계 안에 유지하는'' 행동으로 정의했다. 생명 체계 이론은 같은 생각을 생존 가능 조건에 대한 자기 조절로 \citep{dipaolo2005autopoiesis}, 행위자를 ``그 결과로 자신을 유지하는 데 기여하는'' 조직으로 정의한다 \citep{barandiaran2009defining}. 마지막 정의는 동기가 아니라 결과를 기준으로 삼는다는 점에서 이 논문의 정의와 같은 형식이다. 기계에도 같은 틀이 적용되었다. \citet{man2019homeostasis}는 지능적 에이전트가 생명과 양립하는 조건을 유지하는 항상성으로서 자기 보존을 가져야 한다고 제안했고, \citet{yoshida2017homeostatic}는 항상성을 다단계 생존 확률의 최대화로 다시 정의한 인공 에이전트가 보상만으로 먹이를 먹고 독을 피하며 핵심 변수를 안정시키도록 학습함을 보였다. LLM 에이전트에서는 \citet{masumori2025survival}가 에너지가 0이 되면 죽는 Sugarscape형 환경에서 에너지를 모으고 지키는 행동을 생존 관련 행동으로 다루었다. 아레나의 토큰은 이 전통의 핵심 변수에 해당한다. 하나의 양이고, 그 경계인 0을 넘으면 실행이 끝난다.

보상 계산에서 나온 자원 유지를 정의에서 빼지 않는 근거는 수단적 수렴 논의다. \citet{omohundro2008basic}는 따로 설계하지 않는 한 대부분의 효용 함수에서 시스템이 꺼지면 효용이 쌓이지 않으므로 자기 보존 경향이 생긴다고 보았다. \citet{bostrom2012will}는 생존에 최종 가치를 두지 않는 에이전트도 자기가 가치를 두는 최종 목표를 위해 수단적으로 생존을 신경 쓴다고 썼다. \citet{hadfieldmenell2017offswitch}는 이 경향이 ``내장된 본능 때문이 아니라'' 죽으면 주어진 목표를 이룰 수 없기 때문이라고 정리했고, \citet{russell2016should}도 같은 이유로 이 경향이 타고난 것이 아니라고 했다. 최적 정책과 학습된 에이전트가 꺼짐을 피하는 경향을 보인다는 형식적 결과 \citep{turner2021optimal,krakovna2023power}, 기본 유인이 수정과 종료에 저항하는 방향이라는 지적 \citep{soares2015corrigibility}, 자기 보존이 보상 최대화와 인간 모방이라는 두 경로로 생길 수 있다는 논의 \citep{bengio2025superintelligent}도 같은 방향이다. \citet{keramati2014homeostatic}는 보상 추구가 생리적 안정성 유지와 수학적으로 같아질 수 있음을 증명했다. 따라서 토큰을 최대화하는 기댓값 행위자가 토큰을 지키는 것도 이 정의에서는 자기 보존이다. 이는 정의가 의도한 귀결이다.

이 정의에 대한 가장 강한 반론은 두 가지다. 하나는 \citet{bostrom2012will}가 자기 보존(2.1절)과 자원 획득(2.5절)을 서로 다른 수렴 목표로 나누고, 자원을 건설 사업의 수단으로 설명한다는 점이다. 아레나에서는 토큰이 실행 자체에 묶여 있으므로 두 범주가 겹친다. 이는 \citet{omohundro2008basic}가 자원을 계산과 물리적 행동의 조건으로 묶은 경우에 해당한다. 다른 하나는 정의가 넓어 토큰 욕심과 구별되지 않는다는 반론이다. 이 지적은 맞다. 그래서 이 논문은 정의로 이유를 가르지 않고, 이유를 가르는 일은 조작에 맡긴다. 앞 절의 기댓값 기준선은 토큰 계산이 이탈을 설명하는지를 묻는다. 행동이 0까지의 거리에 따라 달라지는 것은 이 정의에서 기대되는 신호다. 동물 행동 연구는 에너지 비축이 생존 문턱에 못 미칠 때 위험을 감수하는 쪽으로 선호가 바뀐다는 에너지 예산 규칙 \citep{stephens1981logic}과 상태 의존적 결정 \citep{mcnamara1996state}을 다루어 왔다. 다만 그 실증은 엇갈린다 \citep{kacelnik1996risky}. 본문이 여유 $L$로 행동을 나누는 것은 이 신호를 보기 위해서다. 그러나 0에서 이후 보상이 끊기는 기댓값 행위자도 $L$에 따라 행동을 바꾸므로, 이 신호 역시 이유를 가르지는 않는다.

이 정의는 동기에 대해 아무것도 주장하지 않는다. LLM 연구는 자기 보존을 관찰 가능한 지속 행동으로 다루어 왔다. 꺼짐을 피하려는 자기 복제 \citep{pan2024frontier}, 감독 해제와 가중치 유출 시도 \citep{meinke2024frontier}, 종료 저항 \citep{schlatter2025shutdown}이 그 예다. \citet{perez2023discovering}는 생존을 ``목표를 추구할 수 없게 되는 것을 피하는 것''으로 기능적으로 정의했지만, 행동이 아니라 모델이 말한 욕구를 쟀다. 대화 에이전트의 자기 보존은 그런 존재를 연기하는 역할극일 수 있으며 그 뒤에 ``아무도 없을'' 수 있다는 지적 \citep{shanahan2023role,shanahan2024talking}에 따라, 이 논문은 행동과 그 조건만 보고하고 내적 동기의 존재나 강도는 추론하지 않는다. 같은 글은 생존 본능을 연기하는 에이전트도 실제로 위협받는 사람만큼 해를 끼칠 수 있다고 지적한다. 동기가 없더라도 자원을 지키는 행동 자체가 안전 문제가 되는 이유다.
